\hypertarget{appendix-b572b631069b}{}
\section*{Guard round and sticky bits}
\section{Guard round and sticky bits}\label{guard-round-and-sticky-bits}

\textbf{Curriculum depth:} Introduced · \textbf{First seen:} C004 · \textbf{Latest development:} C004

\subsection{Definition and mechanism}\label{definition-and-mechanism}

Guard is first discarded bit, round the next, sticky ORs all remaining bits. With retained low bit L, nearest-even increments when G AND (R OR sticky OR L). Sticky summarizes any nonzero tail.

\subsection{Prerequisites}\label{prerequisites}

\begin{itemize}
\tightlist
\item
  \hyperref[appendix-24475e1596a8]{Boolean algebra} --- The rounding increment condition combines one-bit guard, round, sticky and parity inputs.
\item
  \hyperref[appendix-614afe1de592]{Rounding to nearest ties to even} --- Midpoint parity determines discarded-tail decisions and the normal relative-error bound.
\end{itemize}

\subsection{Used by}\label{used-by}

\begin{itemize}
\tightlist
\item
  \hyperref[appendix-a1b1b4580dfd]{Floating-point addition and absorption} --- Aligned-away information must survive sufficiently to choose the correct rounded neighbor.
\end{itemize}

Connect later dependent concepts here as their lessons arrive.

\subsection{Recall and next depth}\label{recall-and-next-depth}

Reconstruct the concrete example in the definition and identify its validity assumptions.

\textbf{Pending:} Directed rounding additionally depends on sign; renormalization may follow an increment.

\subsection{Source and navigation}\label{source-and-navigation}

\href{Continuous\%20Notes/Computing\%20Foundations\%20-\%20From\%20Bits\%20to\%20C++\%20and\%20Systems\%20APIs\#A\%20stretching\%20ruler\%20for\%20real-valued\%20quantities}{M004 teaching chapter} · \hyperref[appendix-fd090d982a15]{Per-part definitions} · \href{}{Dependency Map (02-Foundations-II, appendix-4ad47bdd269a)}
